\documentclass[10pt,twocolumn]{article}

% ---- packages (standard TeX Live, compila con pdflatex) ----
\usepackage[T1]{fontenc}
\usepackage[utf8]{inputenc}
\usepackage{mathptmx}
\usepackage[margin=2cm]{geometry}
\usepackage{amsmath,amssymb}
\usepackage{booktabs}
\usepackage{titlesec}
\usepackage{fancyhdr}
\usepackage{enumitem}
\usepackage{siunitx}

\titleformat{\section}[block]{\centering\scshape\large}{\thesection.}{0.5em}{}
\titleformat{\subsection}[block]{\bfseries}{\thesubsection}{0.5em}{}
\setlength{\columnsep}{0.8cm}

\pagestyle{fancy}
\fancyhf{}
\fancyhead[R]{\small\itshape Predicting Formula 1 podiums with machine learning}
\renewcommand{\headrulewidth}{0pt}

\title{\bfseries Predicting Formula 1 Podiums with Machine Learning:\\
\large Can a Model Beat the Qualifying Grid?}
\author{%
  Author Name\\ Liceo Scientifico, Italy
}
\date{}

\begin{document}
\twocolumn[%
  \begin{@twocolumnfalse}
    \maketitle
    \begin{center}\textbf{Abstract}\end{center}
    \noindent
    Predicting the outcome of a Formula~1 race is hard: twenty cars, a dominant
    machine factor, and chaotic events (safety cars, weather, reliability). We
    ask a focused question: can a machine learning model predict which drivers
    finish on the \emph{podium} (top~3) better than the single strongest public
    signal, the \emph{qualifying grid}? Using free Ergast data (2010--2026), we
    build a leakage-free feature pipeline (driver and constructor form, track
    history, reliability, grid position) and train a calibrated gradient-boosting
    classifier of $P(\text{podium})$. The model is strong in absolute terms
    (\textbf{AUC $\approx 0.95$}, Brier $0.076$) and well calibrated, but it beats
    the grid-only baseline on the podium-ranking metric (precision@3) by only
    \textbf{$+1.4$ points} over a full season ($0.764$ vs $0.750$). We conclude
    that qualifying already encodes most of the predictive signal---car pace,
    driver form and track---and that machine learning adds a small but real
    improvement on top.
    \vspace{0.6em}
  \end{@twocolumnfalse}
]

% ---------------------------------------------------------------
\begin{center}\scshape\large Indices\end{center}
\begin{enumerate}[leftmargin=1.4em,itemsep=2pt]
  \item[0.] \textbf{Abstract:} summary of problem, method and findings.
  \item[1.] \textbf{Introduction:} why podium prediction, and why the grid is a
            hard baseline.
  \item[2.] \textbf{Related work:} motorsport data and the role of qualifying.
  \item[3.] \textbf{Methodology:} data, leakage-free features, model, calibration
            and evaluation.
  \item[4.] \textbf{Results:} AUC, calibration, precision@3 vs grid, feature
            importance.
  \item[5.] \textbf{Discussion:} car vs driver, limitations.
  \item[6.] \textbf{Conclusion.}
  \item[7.] \textbf{References.}
\end{enumerate}

% ---------------------------------------------------------------
\section{Introduction}
Formula~1 is dominated by the car: a top constructor's driver is far more likely
to reach the podium than a talented driver in a slow car. The result of
qualifying---the starting grid---compresses car pace, driver form and track
suitability into a single number, and is therefore an extremely strong predictor
of the race result. Any useful model must \emph{beat the grid}, not merely
correlate with it.

We focus on a clean, well-posed target: the probability that a driver finishes
on the podium (top~3). This turns a messy 20-way ranking problem into a
calibrated binary classification, and admits a natural ranking metric: of the
three drivers with the highest predicted podium probability, how many actually
make the podium (precision@3). Our research question is whether a machine
learning model can exceed the grid baseline on this metric.

% ---------------------------------------------------------------
\section{Related Work}
Public motorsport data are available through the Ergast API (and its community
mirror jolpica-f1), covering results, grids and lap times since 1950. Analyses of
Formula~1 consistently find that grid position and constructor strength are the
dominant explanatory variables, while race-day chaos (safety cars, weather,
mechanical failures) imposes an irreducible noise floor on any predictor.

% ---------------------------------------------------------------
\section{Methodology}

\subsection{Target}
For each (driver, race) we define the binary label
\begin{equation}
  y = \mathbb{1}[\text{finishing position} \le 3].
\end{equation}
The base rate is $3/20 = 0.15$. We model $P(y=1)$ per driver per race.

\subsection{Data}
We collect race results from the Ergast/jolpica API for seasons 2010--2026:
about \num{7000} (driver, race) records with finishing position, starting grid,
constructor, status (finished / retired) and points. Records are ordered
chronologically by race date.

\subsection{Leakage-free features}
\textbf{Golden rule:} every feature uses only information from \emph{previous}
races; the state is updated only after each race is processed. Features:
\begin{itemize}[leftmargin=1.2em,itemsep=1pt]
  \item \textbf{Grid position} (known pre-race, after qualifying).
  \item \textbf{Driver form}: average finishing position over the last 5 races.
  \item \textbf{Driver points form}: average points over the last 5 races.
  \item \textbf{Driver podium rate}: fraction of podiums in the last 10 races.
  \item \textbf{Driver reliability}: DNF rate over the last 10 races.
  \item \textbf{Constructor form}: average finish of the team (last 12 results),
        a proxy for car pace.
  \item \textbf{Constructor podium rate}.
  \item \textbf{Driver track history}: average finish at this circuit.
  \item \textbf{Round} in the season.
\end{itemize}
Neutral priors are used until enough history accumulates (e.g.\ mid-field finish
$10.5$, podium prior $0.15$).

\subsection{Model and calibration}
We train an XGBoost binary classifier with a strictly temporal split (older
seasons for training, the previous season for validation, the most recent season
held out for testing). Probabilities are calibrated with isotonic regression on
the validation set, so that a predicted ``$70\%$ podium'' truly corresponds to a
$70\%$ frequency. Probability quality is measured by log-loss and the Brier
score.

\subsection{Evaluation}
Because podiums are rare (15\%), raw accuracy is misleading; we report the area
under the ROC curve (AUC) and the Brier score for probability quality, and a
race-level ranking metric:
\begin{equation}
  \text{precision@3} =
  \frac{1}{3R}\sum_{r=1}^{R}\;\sum_{d \in \text{top-3}(r)} y_{d,r},
\end{equation}
where for each of the $R$ races we take the three drivers with the highest
predicted probability. The baseline ranks drivers by grid position
(front of the grid $=$ most likely podium).

% ---------------------------------------------------------------
\section{Results}

\subsection{Probability quality}
The calibrated model discriminates podium from non-podium very well
(Table~\ref{tab:f1q}): an AUC around $0.95$ on a full held-out season.

\begin{table}[t]
\centering
\caption{Probability quality on a held-out season (24 races).}
\label{tab:f1q}
\begin{tabular}{lccc}
\toprule
Metric & AUC & Brier & LogLoss \\
\midrule
Calibrated $P(\text{podium})$ & 0.95 & 0.076 & 0.22 \\
\bottomrule
\end{tabular}
\end{table}

\subsection{Beating the grid}
The decisive comparison is precision@3 against the grid baseline
(Table~\ref{tab:f1p}). Over a full 24-race season the model beats the grid by a
small but real margin; on a tiny 5-race sample the difference is within noise.

\begin{table}[t]
\centering
\caption{Precision@3: of the 3 predicted favourites, how many reach the podium.}
\label{tab:f1p}
\begin{tabular}{lccc}
\toprule
Test set & Model & Grid baseline & $\Delta$ \\
\midrule
Full season (24 races) & \textbf{0.764} & 0.750 & $+1.4$ pt \\
Early season (5 races) & 0.533 & 0.600 & $-6.7$ pt$^\dagger$ \\
\bottomrule
\end{tabular}\\[2pt]
{\footnotesize $^\dagger$ Small-sample noise; not significant.}
\end{table}

\subsection{Feature importance}
The dominant features confirm the intuition: \emph{grid position} is by far the
most important, followed by \emph{driver podium rate} and \emph{constructor}
strength (car pace). Driver track history and reliability contribute marginally.

% ---------------------------------------------------------------
\section{Discussion}
The car (constructor) and the starting grid dominate Formula~1 outcomes, and
qualifying already aggregates car pace, driver form and track suitability into a
single number. A machine learning model that adds rolling form, reliability and
track history improves podium ranking only slightly over the grid alone. This is
consistent with the view that the grid is a near-optimal, low-variance predictor,
and that the remaining variability (safety cars, weather, failures) is largely
irreducible from pre-race data.

\textbf{Limitations.} The most recent season provides only a few races, so its
precision@3 is noisy; a full season is needed for a stable estimate. We did not
evaluate against betting odds (no free historical podium-market odds), so we
measure predictive improvement over the grid, not economic edge. Lap-by-lap
telemetry (e.g.\ FastF1) and weather data could capture in-race dynamics we
ignore.

% ---------------------------------------------------------------
\section{Conclusion}
A calibrated machine learning model predicts Formula~1 podiums very well in
absolute terms (AUC $\approx 0.95$) and improves on the qualifying-grid baseline
by a small but real margin ($+1.4$ precision@3 points over a full season). The
result quantifies how strong the grid already is as a predictor and how much
additional signal driver/constructor form contributes. Future work includes
lap-level telemetry, weather, and extending the binary podium target to full
finishing-order modelling.

% ---------------------------------------------------------------
\section{References}
\begin{enumerate}[leftmargin=1.6em,itemsep=1pt]
  \item Ergast Developer API, \textit{Formula 1 results database}; community
        mirror \textit{jolpica-f1}.
  \item T.~Chen, C.~Guestrin, ``XGBoost: A Scalable Tree Boosting System,''
        \textit{KDD}, 2016.
  \item B.~Zadrozny, C.~Elkan, ``Transforming classifier scores into accurate
        multiclass probability estimates'' (isotonic calibration), \textit{KDD},
        2002.
  \item J.~Davis, M.~Goadrich, ``The relationship between Precision-Recall and
        ROC curves,'' \textit{ICML}, 2006.
\end{enumerate}

\end{document}
